% AI工具使用说明（可独立编译）
% 也可直接粘贴到论文.tex 中替换附录B部分

\documentclass[12pt,a4paper]{ctexart}
\usepackage[top=2.5cm, bottom=2.5cm, left=2.8cm, right=2.8cm]{geometry}
\usepackage{enumitem}
\usepackage{titlesec}
\usepackage{xeCJK}
\usepackage{xcolor}
\usepackage{fancyhdr}
\pagestyle{fancy}
\renewcommand{\headrulewidth}{0pt}
\lhead{} \chead{} \rhead{}
\cfoot{\thepage}

\setCJKmainfont{SimSun}[BoldFont=SimHei]
\setCJKsansfont{SimHei}[AutoFakeBold=3]

\renewcommand{\thesection}{\chinese{section}、}
\newcommand{\bluesectitle}[1]{%
  \colorbox{blue!50!gray!50}{\parbox{\dimexpr\textwidth-2\fboxsep\relax}{\raggedright\heiti\zihao{-3}\bfseries\thesection#1}}}
\titleformat{\section}[block]
  {}
  {}
  {0pt}
  {\bluesectitle}

\begin{document}

\begin{center}
{\zihao{-2}\heiti\bfseries AI 工具使用说明}
\end{center}

\section{AI 工具名称和版本}

DeepSeek V4 Pro、通义千问 3.7 Max

\section{具体使用目的与环节}

\noindent\textbf{\textbullet~DeepSeek V4 Pro}

\begin{enumerate}[label=(\arabic*)]
  \item 文献搜集环节：生成与供应商评价与原材料订购相关的文献中英文检索关键词，梳理CRITIC-TOPSIS综合评价、0-1整数规划、ARIMA时间序列等方向的文献观点，搭建问题一到问题二再到问题三的递进求解框架。
  \item 模型推导环节：整理CRITIC客观赋权与冲突性计算的推导依据、TOPSIS综合排名的数学流程、0-1整数规划约束条件设计逻辑、ARIMA模型定阶与偏差序列构建步骤，对比各阶段算法优劣。
  \item 代码编写环节：生成CRITIC-TOPSIS综合评价、0-1整数规划PuLP求解、ARIMA偏差建模与蒙特卡洛模拟的Python代码骨架，排查求解器参数配置与ARIMA平稳性检验等报错，精简冗余代码。
\end{enumerate}

\noindent\rule{\linewidth}{0.4pt}

\noindent\textbf{\textbullet~通义千问 3.7 Max}

\begin{enumerate}[label=(\arabic*)]
  \item 参考文献排版环节：按GB/T 7714规范将文献信息转换成LaTeX \verb|\bibitem|代码，区分图书、期刊、会议文献类型标识，统一外文作者缩写与标点。
  \item 正文润色环节：优化口语化表述，统一CRITIC权重、0-1整数规划、ARIMA偏差序列等专业术语，结合仿真数据产出图表描述与结果分析初稿。
  \item 格式规整环节：提供竞赛论文排版规范，输出AI使用说明与代码附录的写作模板，规范摘要与结论的行文结构，调整附录表格格式。
\end{enumerate}

\section{关键交互记录}

\noindent\textbf{\textbullet~DeepSeek V4 Pro}

\textbf{Query1：}帮我们梳理供应商评价与原材料订购问题的分析思路，中英文检索词都给一下，再说下问题分析框架怎么写。

\textbf{Output：}关键词：CRITIC-TOPSIS、0-1整数规划、ARIMA偏差预测；分析框架：从供应商综合评价到订购方案优化再到供货预测递进展开。

\textbf{Query2：}推导一下CRITIC客观赋权与冲突性计算的依据，再推导0-1整数规划约束条件和ARIMA偏差序列构建，对比各算法优劣。

\textbf{Output：}CRITIC对比强度与冲突性推导、TOPSIS贴近度数学表述、0-1整数规划产能与频次约束设计、ARIMA定阶与偏差外推步骤；表格对比各算法优劣与局限。

\textbf{Query3：}写个Python代码框架，包括CRITIC-TOPSIS综合评价、0-1整数规划PuLP求解、ARIMA偏差建模与蒙特卡洛模拟，按模块拆分加注释。

\textbf{Output：}分模块代码骨架：特征提取与预处理、CRITIC赋权、TOPSIS排序、0-1规划求解、ARIMA建模、蒙特卡洛模拟；列举求解器配置与ADF检验失效的排查方案。

\noindent\rule{\linewidth}{0.4pt}

\noindent\textbf{\textbullet~通义千问 3.7 Max}

\textbf{Query1：}帮我按GB/T 7714规范把文献信息转成LaTeX参考文献代码，不确定的信息标出来别编。

\textbf{Output：}逐条输出标准化\verb|\bibitem|代码，区分期刊与专著标识；存疑信息单独标注提示人工核验。

\textbf{Query2：}润色一下论文各章节文字，统一专业术语，结合仿真数据写图表分析段落。

\textbf{Output：}改写口语化表述，统一专业词汇；结合仿真数值写图表描述，删减模板套话。

\textbf{Query3：}给竞赛论文排版规范，生成AI使用说明和代码附录的写作模板，调整附录表格格式。

\textbf{Output：}论文标题层级与排版要求；AI使用说明与代码附录写作模板；附录表格多行合并单元格模板代码。

\section{采纳和人工修改情况}

\noindent\textbf{\textbullet~DeepSeek V4 Pro}

\textbf{Query1：}采纳检索关键词与综述写作框架；人工筛选贴合课题的内容，剔除无关解读，自主搭建递进逻辑；所有文献由团队手动核对原始出版信息，修正AI错误年份与页码。

\textbf{Query2：}采纳CRITIC权重推导步骤与算法优劣对比框架；人工补充402家供应商数据特征与55家筛选的业务约束等专属条件，代入全部数据验算公式，修正冲突性计算偏差，自主确定CRITIC与熵权法稳健性验证策略。

\textbf{Query3：}采纳分模块代码框架；人工适配402家供应商数据改写代码，逐行调试修复PuLP求解器配置与ARIMA平稳性检验报错，调整订购方案参数与蒙特卡洛模拟次数等超参数，新增CRITIC冲突性计算与偏差序列自定义函数，独立运行仿真导出全部配图。

\noindent\rule{\linewidth}{0.4pt}

\noindent\textbf{\textbullet~通义千问 3.7 Max}

\textbf{Query1：}采纳LaTeX参考文献格式模板；人工对照数据库逐条核验文献元数据，修正AI添加的错误期号与页码。

\textbf{Query2：}采纳专业语句逻辑；人工结合仿真数据补充量化分析，修正脱离工况的片面解读，核心分析与创新点由团队独立撰写。

\textbf{Query3：}采纳附录排版基础模板与表格代码；人工对照官方模板调整行距与编号，核对附录代码片段完整性，调整章节层级适配全文结构。

\end{document}
